\documentclass[11pt,a4paper]{article}
\usepackage[margin=2.4cm]{geometry}
\usepackage{amsmath,amssymb}
\usepackage{booktabs}
\usepackage{siunitx}
\usepackage{graphicx}
\usepackage[colorlinks=true,linkcolor=blue,citecolor=blue]{hyperref}

\title{Vibronic analysis of the photoexcited states of Fe(acac)$_3$:\\
Franck--Condon activity and Duschinsky mixing}
\author{}
\date{}

\newcommand{\wn}{\ensuremath{\mathrm{cm^{-1}}}}

\begin{document}
\maketitle

\section{Computational methods}

All calculations were carried out with ORCA~6 on iron(III) acetylacetonate,
Fe(acac)$_3$ (43 atoms), in its high-spin sextet ground state ($S=5/2$).
Geometries, energies and gradients were obtained with the TPSSh global hybrid
meta-GGA functional, chosen because it reproduces the experimental absorption
spectrum and NEVPT2 reference excitation energies for this complex better than
the alternatives tested. A mixed basis was used throughout, def2-TZVP on Fe and
O and def2-SVP on C and H, with the def2/J auxiliary basis and the RIJCOSX
approximation. Aqueous solvation was described by the SMD continuum model,
matching the 20\%~ethanol/80\%~water solvent of the experiment; this is
essential rather than cosmetic, because a ligand-to-metal charge-transfer
(LMCT) state redistributes charge and is therefore strongly solvatochromic in
a medium of this polarity, whereas a ligand-centred $\pi\to\pi^{*}$ state is
not. The relative ordering of these states is consequently solvation-controlled.

The ground-state structure was optimised without symmetry constraints and
without any fixed internal coordinate, and confirmed as a true minimum: all
123 vibrational modes are real. Excited states were obtained from
TDA-TD-DFT with 80 roots. Two roots were studied in detail, selected to match
the experimental excitation windows: root~15 at \SI{2.918}{\electronvolt}
(\SI{424.8}{\nano\metre}) and root~61 at \SI{4.475}{\electronvolt}
(\SI{277.0}{\nano\metre}). Both are dominated by $\beta\to\beta$ excitations
into the vacant Fe(III) $\beta$-$d$ manifold, i.e.\ LMCT in character, and both
retain sextet spin ($\langle S^{2}\rangle = 8.80$ and $8.96$ against $8.75$ for
a pure sextet). Excited-state geometries were optimised with root following,
in stages, with the root index re-identified at each stage to account for
reordering along the path; harmonic frequencies were then computed numerically
at the optimised geometries.

\subsection*{Franck--Condon activity from the excited-state gradient}

The primary analysis avoids any dependence on excited-state optimisation. The
excited-state Cartesian gradient $\mathbf{g}$, evaluated at the \emph{ground-state}
geometry, is projected onto the ground-state normal modes
$\mathbf{L}$ (Cartesian columns normalised such that $\sum_i m_i L_{ik}^2 = 1$),
\begin{equation}
g_k = \sum_i L_{ik}\, g_i ,
\end{equation}
in which the mass weighting cancels identically. Within the
vertical-gradient (independent-mode displaced-harmonic-oscillator)
approximation the dimensionless displacement, Huang--Rhys factor and
mode-resolved reorganisation energy follow as
\begin{equation}
\Delta_k = -\frac{g_k}{\omega_k^{3/2}}, \qquad
S_k = \tfrac{1}{2}\Delta_k^{2}, \qquad
\lambda_k = S_k\,\hbar\omega_k ,
\end{equation}
and the resonance-Raman intensity of the fundamental in the short-time
approximation is $I_k \propto \omega_k^{2}\Delta_k^{2}$. Completeness of the
projection was verified in the mass-weighted metric: the residual gradient
unaccounted for by the vibrational modes is 0.16\% and 0.21\% for the two
states, so the decomposition is essentially exact.

\subsection*{Duschinsky analysis}

To compare vibrations between two structures the Duschinsky matrix
\begin{equation}
J_{ij} = \sum_a m_a\, \mathbf{L}^{\mathrm{GS}}_{a,i}\cdot
          \mathcal{R}\,\mathbf{L}^{\mathrm{ES}}_{a,j}
\end{equation}
was constructed, where $\mathcal{R}$ is the rotation that superposes the
excited-state geometry onto the ground-state one, applied to the excited-state
mode vectors so that modes and geometry remain in a common frame. For each
excited-state mode $j$ we report the best-matching ground-state mode, the
normalised overlap $J_{ij}^2/\sum_i J_{ij}^2$, and the participation ratio
\begin{equation}
P_j = \frac{\left(\sum_i J_{ij}^{2}\right)^{2}}{\sum_i J_{ij}^{4}},
\end{equation}
which counts how many ground-state modes are required to build the
excited-state mode: $P_j = 1$ means the mode is preserved intact.

\subsection*{The quartet manifold}

Because relaxation in high-spin Fe(III) may proceed through spin-flipped
states, the quartet manifold was computed at the same ground-state geometry
and the same level of theory using spin-flip TD-DFT (\texttt{SF true}) with
140 roots. This yields 140 quartet states ($\langle S^{2}\rangle \approx 4.0$
against 3.75 for a pure quartet) spanning 1.395--6.632~eV. The 80 roots of the
ordinary TDA-TD-DFT calculation are all sextets, so the two manifolds are
complementary and were combined to give the total density of electronic states
below each excitation.

All analyses were performed in Beavyr; the complete per-mode output is
provided in \texttt{vibronic\_analysis.txt}.

\section{Results and discussion}

\subsection{Ground state and mode characters}

The optimised sextet ground state is a true minimum with 123 real modes. Its
vibrations relevant to the FSRS window were characterised by the per-element
share of the mass-weighted amplitude (Table~\ref{tab:char}). The dense cluster
of twelve modes at 1442--1446~\wn{} is 79\% hydrogen in character and is
therefore methyl deformation; the conjugated C=C/C=O ring stretch lies considerably higher, at
1559.5~\wn{} (C~73\%, O~18\%). The triplet near 811~\wn{} is 87\% hydrogen,
a C--H wag, and is a distinct vibration from $\nu_8$. The sextet of modes at
948--952~\wn{} (C~58\%, H~25\%, O~17\%) corresponds to $\nu_8$, a methyl
rock combined with C--C--C skeletal deformation.

\begin{table}[htbp]
\centering
\caption{Character of the ground-state modes relevant to the FSRS
experiment, as the per-element share of mass-weighted amplitude.}
\label{tab:char}
\begin{tabular}{r l l}
\toprule
Calc.\ / \wn{} & Composition & Assignment \\
\midrule
 277.6 & Fe 67\%, C 18\%, O 9\%  & metal-centred skeletal deformation \\
 435.0 & O 46\%, C 42\%, H 11\%  & Fe--O / chelate ring deformation \\
 661.1 & C 53\%, O 26\%, H 21\%  & ring deformation \\
 811.1 & H 87\%, C 12\%          & C--H wag \\
 948.1 & C 58\%, H 25\%, O 17\%  & $\nu_8$: methyl rock + C--C--C skeleton \\
1031.4 & H 47\%, C 31\%, O 22\%  & $\nu_9$: C--H in-plane bend \\
1295.5 & C 72\%, H 22\%, O 6\%   & C--C stretch \\
1443.9 & H 79\%, C 18\%          & methyl deformation \\
1559.5 & C 73\%, O 18\%, H 9\%   & conjugated C=C/C=O ring stretch \\
\bottomrule
\end{tabular}
\end{table}

The harmonic frequencies are not uniformly scaled relative to experiment: the
$\nu_8$ cluster is computed some 70~\wn{} \emph{above} the observed 878/896~\wn{}
doublet, whereas the $\nu_9$-type C--H bending modes at 1026--1040~\wn{} lie
some 50~\wn{} \emph{below} the observed 1085~\wn{}. Since the functional was
selected to reproduce electronic excitation energies rather than vibrational
frequencies, this is accepted; the arguments below rest on frequency
\emph{shifts} and mode compositions rather than on absolute positions.

\subsection{Franck--Condon activity: where the reorganisation energy goes}

Table~\ref{tab:fc} lists the most Franck--Condon-active modes for the two
states. The total reorganisation energies are
$\lambda = 0.315$~eV (2544~\wn{}) for root~15 and $0.589$~eV (4750~\wn{}) for
root~61: root~61 deposits nearly twice as much energy into nuclear motion.

\begin{table}[htbp]
\centering
\caption{The six most resonance-Raman-active ground-state modes for each
excited state, from the vertical-gradient analysis. Intensities are
normalised to the strongest mode of each state.}
\label{tab:fc}
\begin{tabular}{r r r r  r r r}
\toprule
& \multicolumn{3}{c}{Root 15 (424.8~nm)} & \multicolumn{3}{c}{Root 61 (277.0~nm)} \\
\cmidrule(lr){2-4}\cmidrule(lr){5-7}
Mode / \wn{} & $S_k$ & $\lambda_k$/\wn{} & $I_{\mathrm{RR}}$
             & $S_k$ & $\lambda_k$/\wn{} & $I_{\mathrm{RR}}$ \\
\midrule
1559.5 & 0.173 & 270 & 1.000 & 0.110 & 172 & 0.730 \\
 435.0 & 0.804 & 350 & 0.362 & 1.944 & 846 & 1.000 \\
 200.6 & 3.972 & 797 & 0.380 & 4.788 & 961 & 0.524 \\
 175.3 & 0.710 & 124 & 0.052 & 4.223 & 741 & 0.353 \\
 151.6 & 1.000 & 152 & 0.055 & 4.966 & 753 & 0.311 \\
 194.4 & 0.703 & 137 & 0.063 & 1.928 & 375 & 0.198 \\
\midrule
\multicolumn{1}{l}{Total $\lambda$} & \multicolumn{3}{c}{2544 \wn{} = 0.315 eV}
                                    & \multicolumn{3}{c}{4750 \wn{} = 0.589 eV} \\
\bottomrule
\end{tabular}
\end{table}

Two conclusions follow immediately.

First, the Huang--Rhys factors are substantial. Values of
$S_k = 0.17$--$4.97$ are obtained for the most active modes, and summed over
the experimentally monitored windows they amount to
$S \approx 0.011$--$0.036$. With reorganisation energies of several tenths of
an electronvolt, the vertical transition is accompanied by appreciable
structural reorganisation, which contributes to the transient Raman intensities
alongside resonance enhancement.

Second, and more usefully, the reorganisation is concentrated in modes that
the FSRS measurement does not monitor. The three observed bands together carry
only a few percent of the total: the $\nu_8$ window accounts for 2.4\%
(root~15) and 5.5\% (root~61) of the peak resonance-Raman intensity, and the
$\nu_{10}$ window for 18\% and 16\%. The dominant activity lies instead at
150--450~\wn{} and at 1559~\wn{}. The observed bands are thus reporters of the
vibrational bath rather than of the primary distortion coordinate, which is
consistent with their appearing at both excitation wavelengths and with
nearly identical rise behaviour.

\subsection{The two states are distinguished by Fe--O versus ligand distortion}

The clearest difference between the two states is which coordinate the
excited-state gradient drives (Table~\ref{tab:fc}). For root~15 the strongest
activity is the conjugated C=C/C=O ring stretch at 1559.5~\wn{}, a
ligand-centred coordinate. For root~61 the strongest activity is the
Fe--O/chelate deformation at 435.0~\wn{}, with $S_k = 1.94$ against $0.80$ for
root~15, and the low-frequency skeletal modes at 151--200~\wn{} carry
$\sum S_k \approx 15.9$ against $6.4$. Root~61 therefore distorts the
metal--ligand framework two to three times more strongly than root~15.

This is a mode-resolved, quantitative distinction between the two excitation
windows that does not depend on assigning a qualitative label to either state.
It is notable that the experimental Fe--O symmetric stretch is reported near
449~\wn{}, in close agreement with the computed 435.0~\wn{} mode that carries
the largest Franck--Condon activity of the higher-energy state.

\subsection{Charge transfer must lift the degeneracy of the $E$ vibrations}

A central expectation of this work can be stated before any computed result:
\emph{a ligand-to-metal charge-transfer excitation in a tris-chelate must lift
the degeneracy of the doubly degenerate ligand vibrations}. The transient
spectrum should therefore show the near-degenerate vibrational manifolds of the
ground state split apart, and should do so for any charge-transfer excitation
rather than for one wavelength in particular.

The reasoning is as follows. In its ground state Fe(acac)$_3$ has three
equivalent bidentate ligands and approximate $D_3$ symmetry, so the vibrations
of the ligand framework transform as $A_1$, $A_2$ and $E$, the last doubly
degenerate. Vibrations of the three chelate rings accordingly appear as tight
manifolds: the $\nu_8$ region computed here comprises two ligand modes, each
resolving into a non-degenerate and a doubly degenerate component, with a total
width of only 3.4~\wn{} and an $A$--$E$ separation of 1.3~\wn{}
(Section~\ref{sec:exp}).

A charge-transfer excitation removes an electron from this ligand framework and
places it on the metal. The resulting hole must be accommodated by three
equivalent rings, and it is this circumstance that makes the loss of degeneracy
generic rather than incidental. A state of $D_3$ symmetry requires the hole to
remain equally delocalised over all three rings; whenever vibronic coupling to
the ligand-framework modes exceeds the electronic energy that delocalisation
gains, the delocalised arrangement is unstable and the system relaxes by
localising the hole on one ring, in direct analogy to valence localisation in
mixed-valence compounds. The distortion that achieves this is itself a
ligand-framework deformation, and it destroys the equivalence on which the $E$
degeneracy rests. The degenerate pairs must then separate.

Two features of the present system make this outcome likely. First, the
intrinsic splitting within the $\nu_8$ manifold is very small, 1.3~\wn{}
between the non-degenerate and degenerate components, so only a modest
vibronic coupling is needed to overcome it. Second, both charge-transfer
states lie in a dense region of the excited-state manifold, with neighbouring
states 345~\wn{} away for root~15 and 95~\wn{} for root~61
(Section~\ref{sec:cool}); such close-lying states of appropriate symmetry are
the coupling partners through which second-order vibronic interactions operate.
Neither state is itself orbitally degenerate --- root~15 is separated from its
neighbours by 346 and 1046~\wn{} --- so no first-order Jahn--Teller distortion is
required by symmetry; the mechanism is the second-order, pseudo-Jahn--Teller
one described above.

The prediction is specific and testable: every degenerate ligand manifold of
the ground state, not merely the one the experiment happens to reach, should
appear split in the transient spectrum, and the splitting should accompany
charge-transfer excitation generally rather than being peculiar to one
wavelength. Section~\ref{sec:loc} shows that this is what the calculations
give, for all twelve near-degenerate manifolds and for both excited states.

\subsection{Structural and vibrational evidence for hole localisation}
\label{sec:loc}

The anticipated symmetry breaking is confirmed directly, both in the optimised
excited-state geometries and in the vibrational manifolds as a whole.

\subsubsection*{The FeO$_6$ core}

Table~\ref{tab:feo} lists the six Fe--O distances in each state. In the ground
state they span 0.014~\AA{}, the three chelate rings being equivalent to within
the tolerance of an unconstrained optimisation. Both excited states are
markedly distorted: the spread rises to 0.098~\AA{} for root~15 and 0.241~\AA{}
for root~61, factors of seven and seventeen. The pattern for root~15 is the
one expected of a localised hole, two short distances of 3.849 and
3.852~bohr set against four longer ones between 4.009 and 4.034~bohr, that is,
a single ligand distinguished from the remaining two. Root~61 is distorted
further and less regularly, consistent with its larger reorganisation energy.

\begin{table}[htbp]
\centering
\caption{Fe--O distances in the optimised ground and excited states, in bohr,
sorted by length. The ground state is equivalent to within 0.014~\AA{}; both
excited states single out one chelate ring.}
\label{tab:feo}
\begin{tabular}{l cccccc c c}
\toprule
State & \multicolumn{6}{c}{Fe--O / bohr} & Mean & Spread / \AA{} \\
\midrule
Ground state & 3.797 & 3.797 & 3.808 & 3.808 & 3.823 & 3.823 & 3.809 & 0.014 \\
Root 15      & 3.849 & 3.852 & 4.009 & 4.015 & 4.032 & 4.034 & 3.965 & 0.098 \\
Root 61      & 3.793 & 3.860 & 3.912 & 3.976 & 4.230 & 4.248 & 4.003 & 0.241 \\
\bottomrule
\end{tabular}
\end{table}

\subsubsection*{All ligand manifolds widen, not only $\nu_8$}

If the distortion arises from localisation of the charge-transfer hole, it
should lift the near-degeneracy of every ligand-framework manifold, not merely
the one that happens to be observed. Table~\ref{tab:widen} tests this for all
twelve near-degenerate clusters of three or more modes in the
200--1700~\wn{} interval.

\begin{table}[htbp]
\centering
\caption{Widening of every near-degenerate ground-state cluster in the
excited states. Spans in \wn{}; the final columns give the ratio of excited-
to ground-state span.}
\label{tab:widen}
\begin{tabular}{l r rr rr}
\toprule
& & \multicolumn{2}{c}{Root 15} & \multicolumn{2}{c}{Root 61} \\
\cmidrule(lr){3-4}\cmidrule(lr){5-6}
Ground-state cluster / \wn{} & Span & Span & Factor & Span & Factor \\
\midrule
241.4--246.9   &  5.5 &  15.7 &  2.9 &  53.9 &  9.9 \\
400.3--406.1   &  5.8 & 101.3 & 17.5 &  18.1 &  3.1 \\
562.8--564.0   &  1.2 &  78.7 & 67.8 &  55.5 & 47.9 \\
654.6--667.3   & 12.7 &  88.7 &  7.0 & 167.2 & 13.2 \\
811.1--812.5   &  1.4 &  31.8 & 22.8 &  23.2 & 16.6 \\
948.1--951.5 ($\nu_8$) & 3.4 & 56.6 & 16.7 & 56.3 & 16.6 \\
1026.1--1040.2 & 14.1 & 134.7 &  9.6 &  93.4 &  6.6 \\
1193.6--1196.0 &  2.3 &  27.0 & 11.6 &  32.8 & 14.1 \\
1346.4--1348.4 &  2.0 &   6.8 &  3.3 &  80.5 & 39.5 \\
1370.2--1373.9 &  3.7 &  54.0 & 14.6 &  38.1 & 10.3 \\
1442.1--1446.4 &  4.3 &  43.1 & 10.1 &  73.3 & 17.2 \\
1551.0--1551.8 &  0.8 &  98.9 & 118.6 & 14.2 & 17.0 \\
\midrule
\multicolumn{2}{l}{Median factor} & & 13 & & 15 \\
\bottomrule
\end{tabular}
\end{table}

Every cluster widens, in both states, without exception. The median factor is
13 for root~15 and 15 for root~61, and no cluster widens by less than a factor
of three. The effect is therefore a property of the excited state as a whole
rather than of any individual vibration, which is what a global loss of ligand
equivalence requires and what a local perturbation of one coordinate would not
produce.

The $\nu_8$ manifold, the one accessible to the experiment, is unremarkable
within this set: its factors of 16.7 and 16.6 sit close to the medians. Its
value lies in being observable, not in being exceptional.

\subsection{Duschinsky mixing}

The excited-state stationary points carry a small number of imaginary
frequencies, nine for root~15 and six for root~61. These are soft torsional
motions of the methyl groups, each localised on a different hydrogen and lying
below 200~\wn{}: such near-free rotors have curvature shallow enough that a
residual gradient inverts them, and locating a fully converged minimum on an
excited-state surface of this size is demanding. The Duschinsky analysis is
therefore restricted to the 200--1800~\wn{} interval, which excludes them and
which contains every band relevant to the experiment. Because the Duschinsky
matrix is approximately block diagonal in frequency, the mid-frequency block
is unaffected by the character of the torsional region.

Table~\ref{tab:dusch} summarises the Duschinsky analysis over that block. In
both excited states the mean participation ratio is between 4.3 and 4.6,
meaning that a typical excited-state mode is a mixture of four to five
ground-state modes. Only 12 of 77 modes (root~15) and 7 of 78 (root~61) retain
an overlap above 0.7 with a single ground-state parent.

\begin{table}[htbp]
\centering
\caption{Duschinsky mixing between the ground state and each excited state,
over the 200--1800~\wn{} block.}
\label{tab:dusch}
\begin{tabular}{l r r}
\toprule
& Root 15 & Root 61 \\
\midrule
Imaginary modes below 200~\\wn{} (methyl torsions) & 9 & 6 \\\\
Mean participation ratio $\bar{P}$        & 4.32 & 4.55 \\
Modes with overlap $> 0.7$                & 12 of 77 & 7 of 78 \\
$\nu_8$ cluster span, ground state        & \multicolumn{2}{c}{3.4 \wn{} (948.1--951.5)} \\
$\nu_8$ descendant span, excited state    & 56.6 \wn{} (918.7--975.3) & 56.3 \wn{} (904.8--961.1) \\
\bottomrule
\end{tabular}
\end{table}

The most striking single result is the fate of the $\nu_8$ manifold. In the
ground state the three equivalent acetylacetonate ligands give a near-degenerate
sextet of modes spanning only 3.4~\wn{}. In both excited states this manifold
is shattered, spreading over approximately 57~\wn{} --- a seventeen-fold
widening. This realises the symmetry breaking anticipated above:
the charge-transfer hole localises and the equivalence of the three ligands is
lost. That both states show it, to within 1~\wn{} of the same width despite
differing by 1.6~eV in excitation energy and by a factor of two in
reorganisation energy, indicates that the effect follows from the
charge-transfer character itself rather than from the properties of either
individual state. It constitutes a direct, testable prediction: the excited-state $\nu_8$ region
should appear as a substantially broader multiplet than the narrow ground-state
doublet at 878/896~\wn{}.

The extent of the mixing carries a methodological consequence for the
interpretation of transient spectra. Where a typical excited-state mode is a
combination of four to five ground-state modes, ground-state mode labels do not
transfer to the excited state: describing a transient band as a shifted
$\nu_8$ presumes that $\nu_8$ persists across the transition as an
identifiable normal coordinate, whereas here it is distributed over several
excited-state modes. Assignments of transient features to ground-state labels
are therefore best quoted together with the corresponding Duschinsky overlap,
which in this system is typically 0.3--0.5.

\subsection{Comparison with the measured frequencies and splittings}
\label{sec:exp}

Table~\ref{tab:exp} sets the computed frequencies against the bands observed by
FSRS. The metal--ligand coordinate is reproduced well: the symmetric Fe--O
stretch observed near 449~\wn{} corresponds to the computed mode at
435.0~\wn{}, an error of 14~\wn{}. The ligand-centred modes carry larger and
non-uniform errors, the $\nu_8$ cluster lying some 55--70~\wn{} above the
observed doublet while the $\nu_9$-type C--H bending modes lie 45--59~\wn{}
below the observed band. Since the functional was selected to reproduce
electronic excitation energies rather than vibrational frequencies, absolute
positions are not the basis of the arguments made here; frequency
\emph{shifts} and mode compositions are.

\begin{table}[htbp]
\centering
\caption{Observed bands and their computed counterparts. Errors in absolute
position are non-uniform in sign; the Fe--O coordinate is reproduced most
accurately.}
\label{tab:exp}
\begin{tabular}{l r l r}
\toprule
Band & Obs.\ / \wn{} & Calculated / \wn{} & $\Delta$ / \wn{} \\
\midrule
Fe--O symmetric stretch & 449  & 435.0                    & $-14$ \\
$\nu_8$                 & 878, 896 & 948.1--951.5 (six modes) & $+55$ to $+70$ \\
$\nu_9$                 & 1085 & 1026.1--1040.2 (twelve modes) & $-45$ to $-59$ \\
$\nu_{10}$              & 1456 & 1442.1--1446.4 (methyl deformation) & $-10$ to $-14$ \\
                        &      & 1559.5 (conjugated C=C/C=O)         & $+104$ \\
\midrule
$\nu_8$ splitting, ground state & 18 & 1.3 (A--E), 3.4 (full cluster) & --- \\
\bottomrule
\end{tabular}
\end{table}

\subsubsection*{The ground-state $\nu_8$ splitting}

The observed $\nu_8$ doublet at 878 and 896~\wn{} is separated by 18~\wn{}.
The computed cluster resolves into two ligand modes, each giving a
non-degenerate and a doubly degenerate component,
\begin{center}
948.1, 948.1 (E) \quad 949.4 (A) \quad\quad 950.1, 950.3 (E) \quad 951.5 (A),
\end{center}
with an A--E separation of 1.3~\wn{} and a total cluster width of 3.4~\wn{}.
The calculation therefore underestimates the observed splitting by a factor of
five to fourteen.

The origin of this discrepancy is most probably the solvation model. A
continuum such as SMD is spherically averaged and preserves the equivalence of
the three acetylacetonate ligands essentially exactly, leaving only the
intrinsic through-metal coupling to split the A and E components, which is
genuinely small. In a 20\%~ethanol/80\%~water mixture the three ligands occupy
different instantaneous hydrogen-bonding environments, and the resulting
inhomogeneity is absent from an implicit description. A splitting of 18~\wn{}
is of a reasonable magnitude for such an effect. On this reading the observed
doublet reflects solvational inequivalence of the ligands rather than the
intrinsic symmetry splitting of the isolated complex, a distinction that could
be tested by repeating the ground-state calculation with explicit water
molecules hydrogen-bonded to the chelate oxygens.

\subsubsection*{Consequences for the excited-state prediction}

The ground-state and excited-state splittings reported here have different
physical origins, and the shortfall in the former does not propagate to the
latter. The 18~\wn{} observed in the ground state is attributed above to
solvational inhomogeneity, which the model omits. The 57~\wn{} spread predicted
for the excited states arises instead from an intrinsic electronic effect,
the localisation of charge on one of the three ligands, which the model does
describe.

More importantly, the seventeenfold widening between the ground and excited
states is an internal comparison: the same functional, basis, solvation model
and geometry treatment are applied on both sides, so whatever is absent from
the ground-state description is equally absent from the excited-state one and
largely cancels in the ratio. The two contributions are moreover additive in
sign. An excited state subject to both charge localisation and solvational
inhomogeneity should display a multiplet broader than the computed 57~\wn{},
not narrower, so the prediction of a substantially widened $\nu_8$ region in
the transient spectrum is a conservative one.

\subsection{The 941~\wn{} transient band}

The transient spectra show a band at 941~\wn{}, absent from the ground state
and blue-shifted by 63~\wn{} from the ground-state $\nu_8$ position at
878~\wn{}, at both excitation wavelengths. Three readings of this feature can
be assessed against the calculations.

The first, that it is $\nu_8$ itself stiffened, is not reproduced. The
computed $\nu_8$ descendants shift by $-33$ to $+24$~\wn{} in root~15 and
$-43$ to $+11$~\wn{} in root~61; none approaches $+63$~\wn{}.

The second is that the observed feature corresponds to one of the modes that
the calculation does predict to shift strongly to the blue. The largest such
shifts are $+97.2$~\wn{} (root~15, 1040.2 $\to$ 1137.4~\wn{}) and
$+67.5$~\wn{} (root~61, 1031.4 $\to$ 1098.9~\wn{}). The latter in particular
matches the observed $+63$~\wn{} closely in magnitude. Both parent modes are
C--H in-plane bending vibrations of $\nu_9$ character
(Table~\ref{tab:char}), rather than $\nu_8$. This reading is attractive in
that the calculation reproduces a blue shift of the observed size without
adjustment; it requires, however, that the transient band be reassigned away
from the $\nu_8$ coordinate. It is also consistent with the low Duschinsky
overlaps of these transitions, 0.29 and 0.30, which indicate that the
blue-shifting excited-state modes are strongly mixed and cannot be traced to a
single ground-state parent in any case.

The third possibility is that the ground-state $\nu_8$ splitting is
systematically underestimated. The calculation gives a 3.4~\wn{} span where
experiment resolves 18~\wn{}, a factor of five. If the excited-state span is
underestimated by a comparable factor, the true multiplet would extend well
beyond the computed 57~\wn{} and could readily encompass a component
63~\wn{} above the ground-state origin.

On the present evidence the second reading is the best supported: the
calculation produces excited-state modes blue-shifted by 65--97~\wn{} of the
right magnitude, but their mixed character means the transient band is best
described as a blue-shifted C--H bending/skeletal combination rather than as
$\nu_8$.

\subsection{Vibrational cooling: a density-of-states argument}
\label{sec:cool}

The experiment finds vibrational cooling after 400~nm excitation to be four to
five times faster ($\tau_2 \approx 18$--32~ps) than after 270~nm excitation
($\tau_2 \approx 80$--130~ps). The properties of the two states considered
individually do not account for this ordering, and in isolation would suggest
the opposite: root~61 places more than twice the reorganisation energy into the
low-frequency skeletal modes that couple most efficiently to a solvent bath
($\sum S_k \approx 15.9$ against 6.4) and drives the Fe--O coordinate 2.4
times harder, both of which would favour faster relaxation, while the
Duschinsky mixing metrics differ by only 5\%.

The resolution lies not in the properties of either state but in how many
states lie beneath it. Table~\ref{tab:dos} collects the combined sextet and
quartet manifolds at the Franck--Condon geometry.

\begin{table}[htbp]
\centering
\caption{Electronic states below and around each excitation, from 80 sextet
roots (TDA-TD-DFT) and 140 quartet roots (spin-flip TD-DFT) at the
ground-state geometry.}
\label{tab:dos}
\begin{tabular}{l r r}
\toprule
& Root 15 (2.923~eV) & Root 61 (4.517~eV) \\
\midrule
Sextet states below            & 14 & 60 \\
Quartet states below           & 17 & 59 \\
\textbf{Total states below}    & \textbf{31} & \textbf{119} \\
\midrule
Nearest sextet below / \wn{}   & 1049 & 605 \\
Nearest quartet / \wn{}        & \textbf{250} & 121 \\
Local quartet density / eV$^{-1}$ & 14 & 26 \\
Local sextet density / eV$^{-1}$  & 12 & 24 \\
\midrule
Reorganisation energy / eV     & 0.315 & 0.589 \\
\bottomrule
\end{tabular}
\end{table}

Root~61 has \emph{four times} as many electronic states below it as root~15,
119 against 31. The measured ratio of cooling times is four to five. We
propose that this correspondence is not accidental, and that the measured
$\tau_2$ is not the intrinsic solute--solvent energy-transfer time but the
duration of the electronic cascade, convolved with cooling.

The argument is as follows. Each internal-conversion or intersystem-crossing
step in the cascade converts an electronic energy gap into vibrational quanta.
While the cascade is still running it continuously resupplies vibrational
energy that the solvent is simultaneously removing, so the vibrational
temperature remains elevated for as long as the cascade lasts. For root~15,
2.923~eV is released through some 31 states: the cascade is short, the energy
is delivered promptly, and what is subsequently measured approaches the true
cooling time. For root~61, 4.517~eV is released through some 119 states; the
cascade outlasts the cooling, and the apparent $\tau_2$ is set by the cascade
rather than by the bath. This also explains why the same three modes are
populated with the same rise behaviour at both wavelengths: they are reporters
of a vibrational bath that is being fed from above, consistent with their
carrying only 2--6\% of the reorganisation energy.

The quartet manifold supplies a shortcut, and it is decisive for root~15 in a
way it is not for root~61. Root~15 lies in a sextet-sparse region, 1049~\wn{}
from the nearest same-spin state, but only 250~\wn{} (31~meV) from a quartet.
That gap is comparable to $k_{\mathrm{B}}T$ at room temperature (207~\wn{})
and small against the Fe(III) spin--orbit coupling constant
($\zeta \approx 400$~\wn{}), so the crossing is close to resonant.
Same-spin internal conversion is slow in this region and intersystem crossing
should dominate. Crucially, the quartet manifold extends down to 1.395~eV,
below the lowest sextet at 1.94~eV: a single spin-flip step transfers the
system onto a ladder that reaches further down than the sextet manifold does,
bypassing much of it.

Root~61 is coupled to the quartets even more strongly, the nearest lying only
121~\wn{} away, but crossing gains it nothing: at 4.517~eV it still has 59
quartets beneath it. The shortcut is only useful to a state already near the
bottom of the manifold. Two further observations reinforce this. Root~61's
gaps are roughly half the size of root~15's (605 against 1049~\wn{} for the
nearest sextet), so by the energy-gap law each individual step is faster --- but
each also releases less energy, so more steps are required to dissipate more
total energy, giving a finer and slower staircase. And root~61 drives the
Fe--O coordinate 2.4 times harder, which should promote spin--orbit crossing
more effectively per event; that it nonetheless relaxes more slowly is itself
evidence that the bottleneck is the length of the cascade rather than the rate
of any individual step.

We emphasise the status of this argument. It is a statistical, density-of-states
explanation resting on four mutually consistent quantities --- state counts, gap
sizes, quartet proximity and reorganisation energies --- and its strongest single
element is the agreement between the fourfold ratio of state counts and the
fourfold to fivefold ratio of measured cooling times. It is not a rate
calculation: no spin--orbit matrix elements, Fermi-golden-rule rates or crossing
geometries have been computed. It also rests entirely on TD-DFT state
densities. Individual TD-DFT excitation energies for a complex of this kind
carry uncertainties of a few tenths of an electronvolt, and spin-flip TD-DFT
shows some spin contamination here ($\langle S^{2}\rangle$ up to 4.93 against
3.75). Counting arguments are more forgiving than individual energies, since
errors partly cancel when states are enumerated rather than placed
individually, but the absolute counts remain dependent on the functional and
on the number of roots requested. The proposal should therefore be read as a
plausible and internally consistent mechanism rather than a demonstrated one.
A quantitative test would require spin--orbit couplings
$\langle {}^{6}\Gamma | \hat{H}_{\mathrm{SO}} | {}^{4}\Gamma \rangle$
between the relevant pairs, from which Fermi-golden-rule rates could be
evaluated with the reorganisation energies reported here.

\subsection{Summary}

\begin{enumerate}
\item The sextet ground state is a genuine minimum (123 real modes) at
      TPSSh/SMD(water); the Franck--Condon projection is complete to better
      than 0.21\%.
\item Huang--Rhys factors reach $S_k \approx 5$ and the total reorganisation
      energies are 0.315 and 0.589~eV. Structural reorganisation is
      substantial, not negligible.
\item That reorganisation lies almost entirely outside the monitored bands, at
      150--450~\wn{} and 1559~\wn{}. The observed modes report the vibrational
      bath, not the primary distortion.
\item The two excitation windows are distinguished quantitatively: root~15
      drives the ligand C=C/C=O stretch, root~61 the Fe--O framework two to
      three times more strongly.
\item The 1442--1446~\wn{} cluster is methyl deformation (79\% hydrogen);
      the conjugated C=C/C=O stretch lies at 1559~\wn{}, and the 811~\wn{}
      triplet is a C--H wag distinct from $\nu_8$.
\item Excited-state normal modes are mixtures of four to five ground-state
      modes; ground-state labels should not be transferred to transient bands
      without quoting the Duschinsky overlap.
\item A charge-transfer excitation in a tris-chelate must lift the degeneracy
      of the doubly degenerate $E$ ligand vibrations. A hole shared between
      three equivalent rings localises whenever vibronic coupling exceeds the
      intrinsic splitting of the manifolds, which here is only 1.3~\wn{}, and
      the localising distortion destroys the equivalence the degeneracy rests
      on. Neither state is orbitally degenerate, so the mechanism is
      second-order rather than first-order Jahn--Teller.
\item The prediction is borne out structurally. The six Fe--O distances span
      0.014~\AA{} in the ground state but 0.098 and 0.241~\AA{} in the two
      excited states, root~15 showing two short and four long distances, the
      signature of a hole localised on one chelate ring.
\item It is borne out vibrationally, and globally: all twelve near-degenerate
      ground-state clusters widen in both excited states without exception,
      by a median factor of 13 and 15 and by no less than three. The $\nu_8$
      manifold accessible to experiment widens by 16.7 and 16.6, close to the
      medians --- observable rather than exceptional.
\item The observed 941~\wn{} transient band is best assigned to a
      blue-shifted, strongly mixed C--H bending/skeletal combination rather
      than to a stiffened $\nu_8$.
\item The Fe--O stretch is reproduced to 14~\wn{}. The ground-state $\nu_8$
      splitting is underestimated (1.3--3.4~\wn{} computed against 18~\wn{}
      observed), most probably because a continuum solvation model preserves
      the equivalence of the three ligands that instantaneous hydrogen bonding
      breaks. The predicted excited-state widening is unaffected, being an
      internal comparison at a common level of theory and additive with any
      solvational contribution.
\item Root~61 has four times as many electronic states below it as root~15
      (119 against 31, counting both manifolds), matching the fourfold to
      fivefold ratio of measured cooling times. We propose that $\tau_2$
      measures the duration of the electronic cascade rather than the
      intrinsic cooling time.
\item A near-resonant quartet 250~\wn{} from root~15, in a region where
      same-spin internal conversion is slow, gives the 400~nm excitation a
      spin-flip shortcut onto a manifold reaching below the lowest sextet.
      Root~61 couples to the quartets more strongly still but gains nothing
      by it, having 59 quartets beneath it.
\end{enumerate}

\end{document}
